% Source: https://artofproblemsolving.com/community/c3h1615134_vietas_on_a_forth_degree_polynomial
% Stable identity: aops-topic-1615134-q01
\textbf{1.} Since the leading coefficient is $a$ and the roots are $r_1,\dots,r_4$, the polynomial factors as
\[
ax^4+bx^3+cx^2+dx+e=a(x-r_1)(x-r_2)(x-r_3)(x-r_4).
\]
Expanding the right-hand side,
\[
a\Big(x^4-\big(\textstyle\sum r_i\big)x^3+\big(\sum r_ir_j\big)x^2-\big(\sum r_ir_jr_k\big)x+r_1r_2r_3r_4\Big),
\]
because each term of the expansion chooses $x$ from some factors and $-r_i$ from the others. Comparing coefficients (Vieta's formulas):
\[
\sum r_i=-\dfrac {b}{a},\qquad \sum r_ir_j=\dfrac {c}{a},\qquad \sum r_ir_jr_k=-\dfrac {d}{a},\qquad r_1r_2r_3r_4=\dfrac {e}{a} .
\]
The signs alternate: $-,+,-,+$. For a polynomial of degree $n$ the product of the roots is $(-1)^n\dfrac{\text{constant term}}{\text{leading coefficient}}$, which for $n=4$ is $+\dfrac {e}{a}$.

\textbf{2.} Here $a=2$, $b=-6$, $e=-8$, so the sum of the roots is $-\dfrac{-6}{2}=\boxed{3}$ and the product is $\dfrac{-8}{2}=\boxed{-4}$.
